\documentclass[12pt,a4paper]{report} % Passage en police 12 comme demandé

\usepackage[utf8]{inputenc}
\usepackage[T1]{fontenc}
\usepackage[french,english]{babel} % Ajout de l'anglais

\usepackage{newtxtext,newtxmath} 

\usepackage{geometry}
\geometry{
    paper=a4paper,
    top=2.5cm,
    bottom=2.5cm,
    left=2.5cm,
    right=2.5cm,
    headheight=15pt,  
    footskip=30pt
}

\usepackage{setspace}
\onehalfspacing % Interligne idéal pour les annotations du correcteur

% --- Gestion des tailles de titres ---
\usepackage{titlesec}
\titleformat{\chapter}[block]
  {\normalfont\fontsize{14}{16}\bfseries}{\thechapter.}{0.5em}{}
\titlespacing*{\chapter}{0pt}{25pt}{15pt}

\titleformat{\section}
  {\normalfont\fontsize{12}{14}\bfseries}{\thesection}{1em}{}

\usepackage{fancyhdr}
\pagestyle{fancy}
\fancyhf{} 

\fancyhead[L]{\fontsize{10}{12}\selectfont Computer Science Degree, 3rd year – University of Lille}
\fancyfoot[R]{\thepage}
\fancyfoot[L]{\fontsize{10}{12}\selectfont Stella-Rose MILLE} 

\fancypagestyle{plain}{
    \fancyhf{}
    \fancyhead[L]{\fontsize{10}{12}\selectfont Computer Science Degree, 3rd year – University of Lille}
    \fancyfoot[R]{\thepage}
    \fancyfoot[L]{\fontsize{10}{12}\selectfont Stella-Rose MILLE}
    \renewcommand{\headrulewidth}{0.4pt}
}
\renewcommand{\headrulewidth}{0.4pt}

% --- Figures et Légendes ---
\usepackage{graphicx}
\usepackage{caption}
\captionsetup{font=small, labelfont=bf}

\usepackage[colorlinks=true, linkcolor=black, citecolor=blue, urlcolor=blue]{hyperref}
\usepackage{etoolbox}
\makeatletter
\patchcmd{\chapter}{\if@openright\cleardoublepage\else\clearpage\fi}{}{}{}
\makeatother

\begin{document}
\selectlanguage{english}

% ======================================================================
% EN-TÊTE ET TITRE
% ======================================================================
{\centering
    \vspace*{0.5cm}
    
    % --- Titre du Rapport ---
    {\fontsize{20}{24}\selectfont\bfseries Mid-Term Internship Report}\\[0.4cm]
    {\fontsize{16}{20}\selectfont\bfseries Upgrading the Ume API at INRIA}\\[0.6cm]
    
    \rule{\textwidth}{1pt}\\[0.4cm]

    % --- Intitulé Précis de la Formation ---
    {\fontsize{14}{16}\selectfont\bfseries Computer Science Degree, 3rd year}\\[0.4cm]
    {\fontsize{12}{14}\selectfont Academic Year 2025 -- 2026}\\[0.5cm]
    
    % --- Ligne de séparation ---
    \rule{\textwidth}{1pt}\\[0.4cm]
    
    % --- Bloc Acteurs ---
    {\fontsize{14}{16}\selectfont\scshape Stella-Rose MILLE}\\[1.5cm]
}

% ======================================================================
% CONTENU DU RAPPORT
% ======================================================================

\chapter{Company Overview}
Inria is the national research institute for digital sciences and technology. 
My internship is within the EVREF research team, which specializes in software 
evolution, engineering, and design. They place a particular focus on software 
quality and open-source tools.

\chapter{Preparation and Anticipated Challenges}
Before starting the internship, the main challenge was the programming language. 
The project relies entirely on Pharo, an object-oriented language and environment 
that I did not have the opportunity to learn at university. 
To avoid being completely lost, I learned the basics of the Pharo language 
on my own. Additionally, since the project is open-source, I spent some time 
reading the existing documentation to get an early idea of its architecture. 
This preparation allowed me to overcome my initial stress regarding the 
technical stack and ask the right questions from the very beginning.

\chapter{Expectations}
My main objective during this internship is to upgrade the current Ume API to 
improve its usability for developers. My primary expectations were to build 
confidence in my technical skills and learn to adapt quickly to a new 
technological ecosystem. I also wanted to learn how to write cleaner and 
more maintainable code following industry standards. Finally, I aimed to 
develop my autonomy by implementing features from design to completion.

\chapter{Planned Tasks}
I am working on Ume, a Property-Based Testing (PBT) framework for Pharo. 
Property-Based Testing is an automatic and random test generator. 
Its main advantage is that it can generate a large number of tests to find 
interesting edge cases that could break the code.

I have handled multiple tasks over the past few weeks. 
First, I simplified the campaign definition of Ume to make the configuration 
more intuitive for users. Then, I added a progress bar running on a separate 
thread, which allows users to have real-time feedback on the campaign's 
progress without freezing the environment. I also implemented a shrinking 
method. Shrinking is the capacity to reduce a failing input to find the 
smallest possible case that causes the same bug. It is essential for users 
to understand the origin of a bug more easily so it can be debugged quickly.

\chapter{Review}
To sum up this first period, I am extremely satisfied with the trust and 
autonomy the team has given me. Starting an internship with a completely new 
programming language was intimidating at first. However, I managed to adapt 
quickly and no longer feel lost in the codebase. This experience has 
significantly boosted my self-confidence, proving to me that I am capable of 
learning fast and contributing effectively to a technical team. It strongly 
reinforces my desire to pursue my career path in software development.

\end{document}